\documentclass[a4paper, 12pt]{report}

% Packages
\usepackage{amsmath}
\usepackage{amssymb}
\usepackage[margin=1in]{geometry}
\usepackage{hyperref}
\usepackage{graphicx}
\usepackage{subcaption}
\usepackage{float}
\usepackage{array}
\usepackage{booktabs}
\setcounter{secnumdepth}{3}

\renewcommand{\thesection}{\arabic{section}}
\renewcommand{\thesubsection}{\arabic{section}.\arabic{subsection}}

% Indentation and spacing
\setlength{\parindent}{0pt}

\begin{document}

\section*{Multi-Frequency Analysis}

In this chapter multifrequency simulation campaigns are performed to answer the question that we 
discussed in the first chapter which till date is unsure

\begin{flushleft}
1. Whether number of combinations affect the accuracy of multi-frequency method \\
2. If yes how many combinations in multifrequency method is appropiate and which parameter limits the number of combinations \\
3. Where is the source of error in the multi-frequency method, if it exists \\
\end{flushleft}

Frist, simulation parameters for the multi-frequency analysis are defined based on 
the results obtained from the signal analysis and nonlinearity studies. 
These parameters are then applied to two sets of multi-frequency campaigns, and the results 
are compared with the single-frequency analysis. 
Finally, a framework is formulated for the selection of frequencies and amplitudes for the multi-frequency analysis. 
This chapter is organised as follows:

\begin{flushleft}
\hspace{2cm}1. Simulation Parameters \\
\hspace{2cm}2. Multi-Frequency Analysis - 1 \\
\hspace{2cm}3. Multi-Frequency Analysis - 2 \\
\hspace{2cm}4. Framework for Multi-Frequency Analysis \\
\end{flushleft}

% ============================================================
\section{Simulation Parameters}
% ============================================================

The following criteria govern the selection of frequencies and amplitudes for the 
multi-frequency analysis:

\begin{enumerate}
    \item Frequencies must lie within a frequency bin - derived from the signal analysis.
    \item Amplitude must be less than the threshold amplitude for nonlinearity - derived from the nonlinearity study.
\end{enumerate}

All other simulation parameters are kept the same.\\

Based on these criteria, two sets of multi-frequency analyses are performed. For both simulation 
campaigns, the frequency combinations are identical, but the amplitude selection differs. 
In the first campaign, the individual heave and pitch amplitudes are kept below the nonlinearity 
threshold. In the second campaign, the sum of the heave and pitch motion amplitudes for each 
combination is kept below the threshold. Since the frequency selection is the same for 
both campaigns, the frequency combinations are described once in this section, while the 
respective amplitude selections are described in their corresponding analysis sections.

% ------------------------------------------------------------
\subsection*{Selection of Frequencies and Combinations}
% ------------------------------------------------------------

To evaluate whether increasing the number of combined frequencies affects 
the accuracy of flutter derivative identification, eight frequencies are selected and 
organised into four combination schemes. Combination~1 consists of single-frequency 
simulations and serves as the reference case. Combinations~2, 3, and 4 group these 
frequencies into subsets of two, three, and four, respectively, with the grouping size 
increasing progressively. The frequency selection procedure and the resulting combinations 
are described below. \\

Frequency selection must be made combination wise. In each combination
, the lowest frequency (corresponding to the highest reduced velocity) is 
chosen first, with all other frequencies in the combination selected relative to it. The lowest frequency may be 
chosen independently, but one important condition must be satisfied: its frequency bin spacing 
should be a definite number, so that the dependent frequencies excatly fall in the frequency bins.\\

For both campaigns in this study, lowest frequency is set to $0.125$~Hz, 
giving a corresponding reduced wind speed of $21.86$.

Allowing for 10 oscillations, the total simulation time is:
\begin{equation}
    t_{\text{total}} = \frac{1}{0.125} \times 10 = 80 \text{ s}
\end{equation}

With a time step size of $\Delta t = 0.002$~s, the total number of time steps is $40{,}000$.

The frequency bin spacing for a simulation of 80 seconds are therefore:
\begin{equation}
    \Delta f = \frac{1/\Delta t}{N} = \frac{500}{40000} = 0.0125 \text{ Hz}
\end{equation}

Hence the frequency bins are $[0.0125,\ 0.0250,\ 0.0375,\ 0.0500,\ \ldots]$. \\

\textbf{Note:} The frequency spacing and frequency bins depend on the number of timesteps. Therefore, if the simulation 
is interrupted prematurely, accurate derivative extraction is not guaranteed, as the frequency bins and 
spacing will be altered. \\

Dependent frequencies are selected from this list. An additional recommendation for selection is to avoid selecting 
two frequencies that are adjacent within the frequency bin list. If one frequency does not 
align with a bin due to a shorter-than-expected simulation duration, the result at the 
neighbouring frequency may also be corrupted through spectral leakage. To minimise this risk, 
sufficient spacing between selected frequencies is advised. \\

The frequencies selected for simulation are shown in Table~\ref{tab:frequencies}. 
Frequencies are selected such that each frequency can serve as the primary frequency, 
providing flexibility in combination schemes. The frequencies $0.15$~Hz and $0.25$~Hz are exceptions, 
as they are also needed to cover the reduced velocity range for accurate interpolation of the 
flutter derivative curve.

\begin{table}[htbp]
\centering
\caption{Frequency Parameters for Multi-Frequency Analysis}
\label{tab:frequencies}
\begin{tabular}{@{}lcccccccc@{}}
\toprule
\textbf{Frequency (Hz)} & 0.125 & 0.15 & 0.2 & 0.25 & 0.4 & 0.8 & 1.6 & 3.2 \\
\midrule
Reduced Velocity        & 21.86 & 18.21 & 13.66 & 10.93 & 6.83 & 3.42 & 1.71 & 0.85 \\
Stable Time Period (s)  & 80    & 80    & 50    & 40    & 25   & 12.5 & 6.25 & 3.125 \\
\bottomrule
\end{tabular}
\end{table}

\textbf{Note:} The amplitude is identical for all frequencies; therefore, within each combination, 
only the frequencies are listed and may be matched against the table above.

The combinations are structured as follows. Combination~1 consists of single-frequency simulations, 
which serve as the reference case. The remaining combinations are:

\begin{itemize}
    \item \textbf{Combination 1:} Single frequencies (reference)

    \item \textbf{Combination 2:}
    \begin{itemize}
        \item $[0.125,\ 0.15]$
        \item $[0.2,\ 0.25]$
        \item $[0.4,\ 0.8]$
        \item $[1.6,\ 3.2]$
    \end{itemize}

    \item \textbf{Combination 3:}
    \begin{itemize}
        \item $[0.125,\ 0.15,\ 0.2]$
        \item $[0.25,\ 0.4,\ 0.8]$
        \item $[0.8,\ 1.6,\ 3.2]$
    \end{itemize}

    \item \textbf{Combination 4:}
    \begin{itemize}
        \item $[0.125,\ 0.15,\ 0.2,\ 0.25]$
        \item $[0.4,\ 0.8,\ 1.6,\ 3.2]$
    \end{itemize}
\end{itemize}

% ============================================================
\section{Multi-Frequency Analysis --- 1}
% ============================================================

%TODO: Insert images of multi-frequency input motion signal and aerodynamic force signal,
%      showing separation of components first in the time domain and then in the frequency domain.
%      Do not add frequency-domain plots directly without the time-domain context.

\begin{table}[htbp]
\centering
\renewcommand{\arraystretch}{1.2}
\caption{Frequencies and Amplitude Parameters for Multi-Frequency campaign - 1}
\label{tab:Frequencies and Amplitude Parameters for Multi-Frequency campaign - 1}
\begin{tabular}{|l|cccccccc|}
\hline
\rule{0pt}{12pt}\textbf{Frequency (Hz)} & 0.125 & 0.15 & 0.2 & 0.25 & 0.4 & 0.8 & 1.6 & 3.2 \\[4pt]
\hline
Reduced Velocity        & 21.86 & 18.21 & 13.66 & 10.93 & 6.83 & 3.42 & 1.71 & 0.85 \\
Stable Time Period (s)  & 80    & 80    & 50    & 40    & 25   & 12.5 & 6.25 & 3.125 \\
Heave Amplitude (m)     & 0.0183 & 0.0183 & 0.0183 & 0.0183 & 0.0183 & 0.0183 & 0.0183 & 0.0183 \\
Pitch Amplitude (rad)   & 0.087 & 0.087 & 0.087 & 0.087 & 0.087 & 0.087 & 0.087 & 0.087 \\
\hline
\end{tabular}
\end{table}

\begin{figure}[htbp]
    \centering
    \includegraphics[width=1\textwidth]{Chap_5/Figures/MultiFrequency_Study_1/Combined_Flutter_Derivative.png}
    \caption{Combined Flutter Derivatives --- Campaign 1}
    \label{fig:Combined_Flutter_Derivatives_1}
\end{figure}

\begin{figure}[htbp]
    \centering
    \begin{subfigure}[b]{0.45\textwidth}
        \centering
        \includegraphics[width=\textwidth]{Chap_5/Figures/MultiFrequency_Study_1/Compressed_Displacement_Vs_Force/Heave/NormalizedFy_vs_Combinations_RelativeToFirst.png}
        \caption{}
        \label{fig:MF1_Compiled_Heave_Fy_vs_Combinations}
    \end{subfigure}
    \hfill
    \begin{subfigure}[b]{0.45\textwidth}
        \centering
        \includegraphics[width=\textwidth]{Chap_5/Figures/MultiFrequency_Study_1/Compressed_Displacement_Vs_Force/Pitch/NormalizedFy_vs_Combinations_RelativeToFirst.png}
        \caption{}
        \label{fig:MF1_Compiled_Pitch_Fy_vs_Combinations}
    \end{subfigure}

    \begin{subfigure}[b]{0.45\textwidth}
        \centering
        \includegraphics[width=\textwidth]{Chap_5/Figures/MultiFrequency_Study_1/Compressed_Displacement_Vs_Force/Heave/PhaseFy_vs_Combinations_RelativeToFirst.png}
        \caption{}
        \label{fig:MF1_Compiled_Heave_PhaseFy_vs_Combinations}
    \end{subfigure}
    \hfill
    \begin{subfigure}[b]{0.45\textwidth}
        \centering
        \includegraphics[width=\textwidth]{Chap_5/Figures/MultiFrequency_Study_1/Compressed_Displacement_Vs_Force/Pitch/PhaseFy_vs_Combinations_RelativeToFirst.png}
        \caption{}
        \label{fig:MF1_Compiled_Pitch_PhaseFy_vs_Combinations}
    \end{subfigure}

    \begin{subfigure}[b]{0.45\textwidth}
        \centering
        \includegraphics[width=\textwidth]{Chap_5/Figures/MultiFrequency_Study_1/Compressed_Displacement_Vs_Force/Heave/NormalizedMoment_vs_Combinations_RelativeToFirst.png}
        \caption{}
        \label{fig:MF1_Compiled_Heave_Moment_vs_Combinations}
    \end{subfigure}
    \hfill
    \begin{subfigure}[b]{0.45\textwidth}
        \centering
        \includegraphics[width=\textwidth]{Chap_5/Figures/MultiFrequency_Study_1/Compressed_Displacement_Vs_Force/Pitch/NormalizedMoment_vs_Combinations_RelativeToFirst.png}
        \caption{}
        \label{fig:MF1_Compiled_Pitch_Moment_vs_Combinations}
    \end{subfigure}

    \begin{subfigure}[b]{0.45\textwidth}
        \centering
        \includegraphics[width=\textwidth]{Chap_5/Figures/MultiFrequency_Study_1/Compressed_Displacement_Vs_Force/Heave/PhaseMx_vs_Combinations_RelativeToFirst.png}
        \caption{}
        \label{fig:MF1_Compiled_Heave_PhaseMx_vs_Combinations}
    \end{subfigure}
    \hfill
    \begin{subfigure}[b]{0.45\textwidth}
        \centering
        \includegraphics[width=\textwidth]{Chap_5/Figures/MultiFrequency_Study_1/Compressed_Displacement_Vs_Force/Pitch/PhaseMx_vs_Combinations_RelativeToFirst.png}
        \caption{}
        \label{fig:MF1_Compiled_Pitch_PhaseMx_vs_Combinations}
    \end{subfigure}

    \caption{Compiled Displacement vs.\ Force --- Campaign 1}
    \label{fig:MF1_Compiled_Displacement_Vs_Force}
\end{figure}

% ------------------------------------------------------------
\subsection{Key Findings}
% ------------------------------------------------------------

\begin{itemize}
    \item Flutter derivatives $H_1^*$, $H_3^*$, $A_1^*$, and $A_3^*$ are accurately identified in the 
          multi-frequency forced-motion CFD simulation.

    \item Flutter derivatives $H_2^*$, $H_4^*$, $A_2^*$, and $A_4^*$ exhibit errors in the 
          multi-frequency forced-motion CFD simulation.

    \item The frequency-domain response under pitch excitation contains higher-order harmonic 
          terms, indicating the presence of nonlinearity. This affects the pitch-motion-associated 
          derivatives: $H_2^*$, $H_3^*$, $A_2^*$, and $A_3^*$.

    \item These errors have a significant influence on the predicted critical flutter velocity.
\end{itemize}

% ============================================================
\section{Multi-Frequency Analysis --- 2}
% ============================================================


\begin{table}[htbp]
\centering
\renewcommand{\arraystretch}{1.2}
\caption{Frequencies and Amplitude Parameters for Multi-Frequency campaign - 1}
\label{tab:Frequencies and Amplitude Parameters for Multi-Frequency campaign - 1}
\begin{tabular}{|l|c|c|c|c|c|c|c|c|}
\hline
\rule{0pt}{12pt}\textbf{MF2 Simulation} & 1 & 2 & 3 & 4 & 5 & 6 & 7 & 8 \\[4pt]
\hline
\rule{0pt}{12pt}\textbf{Frequency (Hz)} & 0.125 & 0.15 & 0.2 & 0.25 & 0.4 & 0.8 & 1.6 & 3.2 \\[4pt]
\hline
Reduced Velocity        & 21.86 & 18.21 & 13.66 & 10.93 & 6.83 & 3.42 & 1.71 & 0.85 \\
Stable Time Period (s)  & 80    & 80    & 50    & 40    & 25   & 12.5 & 6.25 & 3.125 \\
Heave Amplitude (m)     & 0.01 & 0.01 & 0.01 & 0.01 & 0.01 & 0.01 & 0.01 & 0.01 \\
Pitch Amplitude (rad)   & 0.054 & 0.054 & 0.054 & 0.054 & 0.054 & 0.054 & 0.054 & 0.054 \\
\hline
\end{tabular}
\end{table}

The prescribed forced motion for the MF2 simulation 9 is shown in Figure


\begin{table}[htbp]
\centering
\renewcommand{\arraystretch}{1.2}
\caption{Frequencies and Amplitude Parameters for Multi-Frequency campaign - 1}
\label{tab:Frequencies and Amplitude Parameters for Multi-Frequency campaign - 1}
\begin{tabular}{|l|cc|cc|cc|cc|}
\hline
\rule{0pt}{12pt}\textbf{MF2 Simulation} & \multicolumn{2}{c|}{9} & \multicolumn{2}{c|}{10} & \multicolumn{2}{c|}{11} & \multicolumn{2}{c|}{12} \\[4pt]
\hline
\rule{0pt}{12pt}\textbf{Frequency (Hz)} & 0.125 & 0.15 & 0.2 & 0.25 & 0.4 & 0.8 & 1.6 & 3.2 \\[4pt]
\hline
Reduced Velocity        & 21.86 & 18.21 & 13.66 & 10.93 & 6.83 & 3.42 & 1.71 & 0.85 \\
Stable Time Period (s)  & 80    & 80    & 50    & 40    & 25   & 12.5 & 6.25 & 3.125 \\
Heave Amplitude (m)     & 0.009 & 0.009 & 0.009 & 0.009 & 0.009 & 0.009 & 0.009 & 0.009 \\
Pitch Amplitude (rad)   & 0.027 & 0.027 & 0.027 & 0.027 & 0.027 & 0.027 & 0.027 & 0.027 \\
\hline
\end{tabular}
\end{table}

The prescirbed forced motion for the MF2 simulation 9 is shown in Figure



\begin{table}[htbp]
\centering
\renewcommand{\arraystretch}{1.2}
\caption{Frequencies and Amplitude Parameters for Multi-Frequency campaign - 1}
\label{tab:Frequencies and Amplitude Parameters for Multi-Frequency campaign - 1}
\begin{tabular}{|l|ccc|ccc|ccc|}
\hline
\rule{0pt}{12pt}\textbf{MF2 Simulation} & \multicolumn{3}{c|}{13} & \multicolumn{3}{c|}{14} & \multicolumn{3}{c|}{15} \\[4pt]
\hline
\rule{0pt}{12pt}\textbf{Frequency (Hz)} & 0.125 & 0.15 & 0.2 & 0.25 & 0.4 & 0.8 & 0.8 & 1.6 & 3.2 \\[4pt]
\hline
Reduced Velocity        & 21.86 & 18.21 & 13.66 & 10.93 & 6.83 & 3.42 & 3.42 & 1.71 & 0.85 \\
Stable Time Period (s)  & 80    & 80    & 50    & 40    & 25   & 12.5 & 12.5 & 6.25 & 3.125 \\
Heave Amplitude (m)     & 0.006 & 0.006 & 0.006 & 0.006 & 0.006 & 0.006 & 0.006 & 0.006 & 0.006 \\
Pitch Amplitude (rad)   & 0.017 & 0.017 & 0.017 & 0.017 & 0.017 & 0.017 & 0.017 & 0.017 & 0.017 \\
\hline
\end{tabular}
\end{table}

The prescribed forced motion for the MF2 simulation 9 is shown in Figure

\begin{table}[htbp]
\centering
\renewcommand{\arraystretch}{1.2}
\caption{Frequencies and Amplitude Parameters for Multi-Frequency campaign - 1}
\label{tab:Frequencies and Amplitude Parameters for Multi-Frequency campaign - 1}
\begin{tabular}{|l|cccc|cccc|}
\hline
\rule{0pt}{12pt}\textbf{MF2 Simulation} & \multicolumn{4}{c|}{16} & \multicolumn{4}{c|}{17} \\[4pt]
\hline
\rule{0pt}{12pt}\textbf{Frequency (Hz)} & 0.125 & 0.15 & 0.2 & 0.25 & 0.4 & 0.8 & 1.6 & 3.2 \\[4pt]
\hline
Reduced Velocity        & 21.86 & 18.21 & 13.66 & 10.93 & 6.83 & 3.42 & 1.71 & 0.85 \\
Stable Time Period (s)  & 80    & 80    & 50    & 40    & 25   & 12.5 & 6.25 & 3.125 \\
Heave Amplitude (m)     & 0.0045 & 0.0045 & 0.0045 & 0.0045 & 0.0045 & 0.0045 & 0.0045 & 0.0045 \\
Pitch Amplitude (rad)   & 0.013 & 0.013 & 0.013 & 0.013 & 0.013 & 0.013 & 0.013 & 0.013 \\
\hline
\end{tabular}
\end{table}

The prescribed forced motion for the MF2 simulation 9 is shown in Figure

A explanation procedure of multifrequency from prescribed forced motion to obtaining flutter derivatives for Multifrequency
method is given below for MF2 Simulation 16 with Heave motion:

\begin{figure}[htbp]
    \centering
    \includegraphics[width=1\textwidth]{Chap_5/Figures/MultiFrequency_Study_2/Procedure_Explanation/2_Individual_Component_ForcedMotion.PNG}
    \caption{Combined Flutter Derivatives --- Campaign 2}
    \label{fig:Combined_Flutter_Derivatives_2}
\end{figure}

\begin{figure}[htbp]
    \centering
    \includegraphics[width=1\textwidth]{Chap_5/Figures/MultiFrequency_Study_2/Procedure_Explanation/2_Superposed_ForcedMotion.png}
    \caption{Combined Flutter Derivatives --- Campaign 2}
    \label{fig:Combined_Flutter_Derivatives_2}
\end{figure}

\begin{figure}[htbp]
    \centering
    \includegraphics[width=1\textwidth]{Chap_5/Figures/MultiFrequency_Study_2/Procedure_Explanation/3_Computed_Aerodynamic_Lift_Force.png}
    \caption{Combined Flutter Derivatives --- Campaign 2}
    \label{fig:Combined_Flutter_Derivatives_2}
\end{figure}

\begin{figure}[htbp]
    \centering
    \includegraphics[width=1\textwidth]{Chap_5/Figures/MultiFrequency_Study_2/Procedure_Explanation/3_DFT_LiftForce_Superposition_Case_4.png}
    \caption{Combined Flutter Derivatives --- Campaign 2}
    \label{fig:Combined_Flutter_Derivatives_2}
\end{figure}

\begin{figure}[htbp]
    \centering
    \includegraphics[width=1\textwidth]{Chap_5/Figures/MultiFrequency_Study_2/Procedure_Explanation/3_Individual_Harmonic_Lift_Force.png}
    \caption{Combined Flutter Derivatives --- Campaign 2}
    \label{fig:Combined_Flutter_Derivatives_2}
\end{figure}


\begin{figure}[htbp]
    \centering
    \includegraphics[width=1\textwidth]{Chap_5/Figures/MultiFrequency_Study_2/Procedure_Explanation/4_Computed_Aerodynamic_Moment.png}
    \caption{Combined Flutter Derivatives --- Campaign 2}
    \label{fig:Combined_Flutter_Derivatives_2}
\end{figure}

\begin{figure}[htbp]
    \centering
    \includegraphics[width=1\textwidth]{Chap_5/Figures/MultiFrequency_Study_2/Procedure_Explanation/4_DFT_Moment_Superposition_Case_4.png}
    \caption{Combined Flutter Derivatives --- Campaign 2}
    \label{fig:Combined_Flutter_Derivatives_2}
\end{figure}

\begin{figure}[htbp]
    \centering
    \includegraphics[width=1\textwidth]{Chap_5/Figures/MultiFrequency_Study_2/Procedure_Explanation/4_Individual_Harmonic_Moment.png}
    \caption{Combined Flutter Derivatives --- Campaign 2}
    \label{fig:Combined_Flutter_Derivatives_2}
\end{figure}



\begin{figure}[htbp]
    \centering
    \includegraphics[width=1\textwidth]{Chap_5/Figures/MultiFrequency_Study_2/Combined_Flutter_Derivative.png}
    \caption{Combined Flutter Derivatives --- Campaign 2}
    \label{fig:Combined_Flutter_Derivatives_2}
\end{figure}

\begin{figure}[htbp]
    \centering
    \includegraphics[width=1\textwidth]{Chap_5/Figures/MultiFrequency_Study_2/ReducedVelocity_Vs_Number_of_Combinations.png}
    \caption{Reduced Velocity vs.\ Number of Combinations}
    \label{fig:ReducedVelocity_Vs_Number_of_Combinations}
\end{figure}

\begin{figure}[htbp]
    \centering
    \includegraphics[width=1\textwidth]{Chap_5/Figures/MultiFrequency_Study_2/Combined_Amplitude_Errors.png}
    \caption{Combined Amplitude Errors}
    \label{fig:Combined_Amplitude_Errors}
\end{figure}

\begin{figure}[htbp]
    \centering
    \includegraphics[width=1\textwidth]{Chap_5/Figures/MultiFrequency_Study_2/Combined_Phase_Errors.png}
    \caption{Combined Phase Errors}
    \label{fig:MF2_Combined_Phase_Errors}
\end{figure}


\begin{figure}[htbp]  
    \centering
    \begin{subfigure}[b]{0.45\textwidth}
        \centering
        \includegraphics[width=\textwidth]{Chap_5/Figures/MultiFrequency_Study_2/Frequency_Domain_Analysis/Heave/DFT_LiftForce_Superposition_Case_1.png}
        \caption{}
        \label{fig:Heave_DFT_Fy_amplitude_case_0}
    \end{subfigure}
    \hfill
    \begin{subfigure}[b]{0.45\textwidth}
        \centering
        \includegraphics[width=\textwidth]{Chap_5/Figures/MultiFrequency_Study_2/Frequency_Domain_Analysis/Heave/DFT_Moment_Superposition_Case_1.png}
        \caption{}
        \label{fig:Heave_DFT_moment_amplitude_case_0}
    \end{subfigure}
    
    \begin{subfigure}[b]{0.45\textwidth}
        \centering
        \includegraphics[width=\textwidth]{Chap_5/Figures/MultiFrequency_Study_2/Frequency_Domain_Analysis/Heave/DFT_LiftForce_Superposition_Case_2.png}
        \caption{}
        \label{fig:Heave_DFT_Fy_amplitude_case_2}
    \end{subfigure}
    \hfill
    \begin{subfigure}[b]{0.45\textwidth}
        \centering
        \includegraphics[width=\textwidth]{Chap_5/Figures/MultiFrequency_Study_2/Frequency_Domain_Analysis/Heave/DFT_Moment_Superposition_Case_2.png}
        \caption{}
        \label{fig:Heave_DFT_moment_amplitude_case_2}
    \end{subfigure}
    
    \begin{subfigure}[b]{0.45\textwidth}
        \centering
        \includegraphics[width=\textwidth]{Chap_5/Figures/MultiFrequency_Study_2/Frequency_Domain_Analysis/Heave/DFT_LiftForce_Superposition_Case_3.png}
        \caption{}
        \label{fig:Heave_DFT_Fy_amplitude_case_4}
    \end{subfigure}
    \hfill
    \begin{subfigure}[b]{0.45\textwidth}
        \centering
        \includegraphics[width=\textwidth]{Chap_5/Figures/MultiFrequency_Study_2/Frequency_Domain_Analysis/Heave/DFT_Moment_Superposition_Case_3.png}
        \caption{}
        \label{fig:Heave_DFT_moment_amplitude_case_4}
    \end{subfigure}
    
    \begin{subfigure}[b]{0.45\textwidth}
        \centering
        \includegraphics[width=\textwidth]{Chap_5/Figures/MultiFrequency_Study_2/Frequency_Domain_Analysis/Heave/DFT_LiftForce_Superposition_Case_4.png}
        \caption{}
        \label{fig:Heave_DFT_Fy_amplitude_case_6}
    \end{subfigure}
    \hfill
    \begin{subfigure}[b]{0.45\textwidth}
        \centering
        \includegraphics[width=\textwidth]{Chap_5/Figures/MultiFrequency_Study_2/Frequency_Domain_Analysis/Heave/DFT_Moment_Superposition_Case_4.png}
        \caption{}
        \label{fig:Heave_DFT_moment_amplitude_case_6}
    \end{subfigure}
    
    \caption{Frequency Domain Analysis Heave}
    \label{fig:Frequency_Domain_Analysis_Heave}
\end{figure}


\begin{figure}[htbp]  
    \centering
    \begin{subfigure}[b]{0.45\textwidth}
        \centering
        \includegraphics[width=\textwidth]{Chap_5/Figures/MultiFrequency_Study_2/Frequency_Domain_Analysis/Pitch/DFT_LiftForce_Superposition_Case_1.png}
        \caption{}
        \label{fig:Heave_DFT_Fy_amplitude_case_0}
    \end{subfigure}
    \hfill
    \begin{subfigure}[b]{0.45\textwidth}
        \centering
        \includegraphics[width=\textwidth]{Chap_5/Figures/MultiFrequency_Study_2/Frequency_Domain_Analysis/Pitch/DFT_Moment_Superposition_Case_1.png}
        \caption{}
        \label{fig:Heave_DFT_moment_amplitude_case_0}
    \end{subfigure}
    
    \begin{subfigure}[b]{0.45\textwidth}
        \centering
        \includegraphics[width=\textwidth]{Chap_5/Figures/MultiFrequency_Study_2/Frequency_Domain_Analysis/Pitch/DFT_LiftForce_Superposition_Case_2.png}
        \caption{}
        \label{fig:Heave_DFT_Fy_amplitude_case_2}
    \end{subfigure}
    \hfill
    \begin{subfigure}[b]{0.45\textwidth}
        \centering
        \includegraphics[width=\textwidth]{Chap_5/Figures/MultiFrequency_Study_2/Frequency_Domain_Analysis/Pitch/DFT_Moment_Superposition_Case_2.png}
        \caption{}
        \label{fig:Heave_DFT_moment_amplitude_case_2}
    \end{subfigure}
    
    \begin{subfigure}[b]{0.45\textwidth}
        \centering
        \includegraphics[width=\textwidth]{Chap_5/Figures/MultiFrequency_Study_2/Frequency_Domain_Analysis/Pitch/DFT_LiftForce_Superposition_Case_3.png}
        \caption{}
        \label{fig:Heave_DFT_Fy_amplitude_case_4}
    \end{subfigure}
    \hfill
    \begin{subfigure}[b]{0.45\textwidth}
        \centering
        \includegraphics[width=\textwidth]{Chap_5/Figures/MultiFrequency_Study_2/Frequency_Domain_Analysis/Pitch/DFT_Moment_Superposition_Case_3.png}
        \caption{}
        \label{fig:Heave_DFT_moment_amplitude_case_4}
    \end{subfigure}
    
    \begin{subfigure}[b]{0.45\textwidth}
        \centering
        \includegraphics[width=\textwidth]{Chap_5/Figures/MultiFrequency_Study_2/Frequency_Domain_Analysis/Pitch/DFT_LiftForce_Superposition_Case_4.png}
        \caption{}
        \label{fig:Heave_DFT_Fy_amplitude_case_6}
    \end{subfigure}
    \hfill
    \begin{subfigure}[b]{0.45\textwidth}
        \centering
        \includegraphics[width=\textwidth]{Chap_5/Figures/MultiFrequency_Study_2/Frequency_Domain_Analysis/Pitch/DFT_Moment_Superposition_Case_4.png}
        \caption{}
        \label{fig:Heave_DFT_moment_amplitude_case_6}
    \end{subfigure}
    
    \caption{Frequency Domain Analysis Heave}
    \label{fig:Frequency_Domain_Analysis_Heave}
\end{figure}


\begin{figure}[htbp]
    \centering
    \begin{subfigure}[b]{0.45\textwidth}
        \centering
        \includegraphics[width=\textwidth]{Chap_5/Figures/MultiFrequency_Study_2/Compressed_Displacement_Vs_Force/Heave/NormalizedFy_vs_Combinations_RelativeToFirst.png}
        \caption{}
        \label{fig:MF2_Compiled_Heave_Fy_vs_Combinations}
    \end{subfigure}
    \hfill
    \begin{subfigure}[b]{0.45\textwidth}
        \centering
        \includegraphics[width=\textwidth]{Chap_5/Figures/MultiFrequency_Study_2/Compressed_Displacement_Vs_Force/Pitch/NormalizedFy_vs_Combinations_RelativeToFirst.png}
        \caption{}
        \label{fig:MF2_Compiled_Pitch_Fy_vs_Combinations}
    \end{subfigure}

    \begin{subfigure}[b]{0.45\textwidth}
        \centering
        \includegraphics[width=\textwidth]{Chap_5/Figures/MultiFrequency_Study_2/Compressed_Displacement_Vs_Force/Heave/PhaseFy_vs_Combinations_RelativeToFirst.png}
        \caption{}
        \label{fig:MF2_Compiled_Heave_PhaseFy_vs_Combinations}
    \end{subfigure}
    \hfill
    \begin{subfigure}[b]{0.45\textwidth}
        \centering
        \includegraphics[width=\textwidth]{Chap_5/Figures/MultiFrequency_Study_2/Compressed_Displacement_Vs_Force/Pitch/PhaseFy_vs_Combinations_RelativeToFirst.png}
        \caption{}
        \label{fig:MF2_Compiled_Pitch_PhaseFy_vs_Combinations}
    \end{subfigure}

    \begin{subfigure}[b]{0.45\textwidth}
        \centering
        \includegraphics[width=\textwidth]{Chap_5/Figures/MultiFrequency_Study_2/Compressed_Displacement_Vs_Force/Heave/NormalizedMoment_vs_Combinations_RelativeToFirst.png}
        \caption{}
        \label{fig:MF2_Compiled_Heave_Moment_vs_Combinations}
    \end{subfigure}
    \hfill
    \begin{subfigure}[b]{0.45\textwidth}
        \centering
        \includegraphics[width=\textwidth]{Chap_5/Figures/MultiFrequency_Study_2/Compressed_Displacement_Vs_Force/Pitch/NormalizedMoment_vs_Combinations_RelativeToFirst.png}
        \caption{}
        \label{fig:MF2_Compiled_Pitch_Moment_vs_Combinations}
    \end{subfigure}

    \begin{subfigure}[b]{0.45\textwidth}
        \centering
        \includegraphics[width=\textwidth]{Chap_5/Figures/MultiFrequency_Study_2/Compressed_Displacement_Vs_Force/Heave/PhaseMx_vs_Combinations_RelativeToFirst.png}
        \caption{}
        \label{fig:MF2_Compiled_Heave_PhaseMx_vs_Combinations}
    \end{subfigure}
    \hfill
    \begin{subfigure}[b]{0.45\textwidth}
        \centering
        \includegraphics[width=\textwidth]{Chap_5/Figures/MultiFrequency_Study_2/Compressed_Displacement_Vs_Force/Pitch/PhaseMx_vs_Combinations_RelativeToFirst.png}
        \caption{}
        \label{fig:MF2_Compiled_Pitch_PhaseMx_vs_Combinations}
    \end{subfigure}

    \caption{Compiled Displacement vs.\ Force --- Campaign 2}
    \label{fig:MF2_Compiled_Displacement_Vs_Force}
\end{figure}

\begin{figure}[htbp]
    \centering
    \includegraphics[width=1\textwidth]{Chap_5/Figures/Simulation_Time.png}
    \caption{Simulation Time Comparison}
    \label{fig:Simulation_Time}
\end{figure}

% ============================================================
\section{Framework for Multi-Frequency Analysis}
% ============================================================

Based on the results obtained above, the following framework is proposed for conducting a 
multi-frequency analysis.

\begin{enumerate}
    \item \textbf{Determine the number of combinations.} The recommended number is 3 or 4.

    \item \textbf{Target the lower frequency range.} Combinations should preferably be formed 
          from the lowest frequencies (i.e., those at the highest reduced velocities).

    \item \textbf{Select one independent frequency per combination.}

    \item \textbf{Derive the dependent frequencies} within each combination such that they 
          fall within the frequency bin of the independent frequency. A detailed procedure 
          for this calculation is given below.

    \item \textbf{Select amplitudes} such that the total amplitude (sum across all frequencies 
          in a combination) remains below a heave ratio of 5\% and a pitch angle of 3°.

    \item \textbf{Post-processing:} Apply the Discrete Fourier Transform (DFT) to the 
          aerodynamic forces and moments to extract individual frequency components. 
          Also inspect the spectrum for signs of nonlinearity.
\end{enumerate}

\vspace{0.5em}
\textbf{Important Notes:}

\begin{itemize}
    \item Because DFT is used, accurate spectral decomposition of the aerodynamic forces is 
          only achievable for frequency components that fall exactly within a frequency bin. 
          Accurate results cannot be obtained for frequencies that fall between bins.

    \item Even if a frequency of interest falls exactly within a frequency bin, spectral leakage 
          from a neighbouring frequency that does \emph{not} fall within a bin can contaminate 
          the extracted amplitude and phase at that frequency.

    \item Therefore, when selecting frequencies, avoid choosing two that are consecutive 
          in the frequency bin list, as sufficient spacing between selected frequencies 
          reduces the impact of spectral leakage.
\end{itemize}

\vspace{0.5em}
\textbf{Selection of Dependent Frequencies:}

\end{document}